\section{Reed--Solomon 编码}
\label{sec:04-Discrete-Mathematics/05-Coding-and-Polynomial-Methods/03}

回到导读里那个条带：6 块数据，3 块校验，9 台机器各存一块。运维手册上写的容错承诺是``任意 3 台同时失效仍可恢复''。注意``任意''这两个字——不是``最多坏 3 台校验机''，也不是``坏的那几台恰好不相邻''，是九台里随便哪三台。

三副本给不出这种承诺，它只知道复制。要让九块存储彼此对称到这个程度，得让每一块都是同一个对象的一次采样，而不是某一块的影子。上一节末尾已经点了这件事：把消息装进一个低次多项式，条带上的九块就是它在九个位置上的取值。

\subsection{把消息当成系数，把码字当成取值}

固定有限域 \(\mathbb F_q\)，在其中挑 \(n\) 个互异的求值点 \(\alpha_1,\ldots,\alpha_n\)，于是必须有 \(n\le q\)。把 \(k\) 个消息符号 \(m_0,\ldots,m_{k-1}\) 当作多项式
\[
f(x)=m_0+m_1x+\cdots+m_{k-1}x^{k-1}
\]
的系数，码字就是
\[
\bigl(f(\alpha_1),\ldots,f(\alpha_n)\bigr).
\]
这便是 \((n,k)\) Reed--Solomon 码，码率 \(k/n\)。消息一共 \(k\) 个自由度，却被摊成 \(n\) 个互相牵制的符号。有些实现采用系统形式，让前 \(k\) 个位置直接等于消息本身，那只是给同一个线性码换了一组消息坐标，构造并没有变。

\begin{center}
\begin{tikzpicture}[x=1cm,y=1cm,font=\small,>=stealth]
  \foreach \i in {0,1,2,3,4,5,6,7,8}{
    \draw[black!50,fill=blue!10] (\i,3.4) rectangle (\i+0.9,4.2);
  }
  \foreach \i in {2,5,8}{
    \draw[black!45,fill=black!12] (\i,3.4) rectangle (\i+0.9,4.2);
    \draw[black!60] (\i+0.18,3.58) -- (\i+0.72,4.02);
    \draw[black!60] (\i+0.18,4.02) -- (\i+0.72,3.58);
  }
  \node[above,font=\footnotesize,black!70] at (4.5,4.3) {9 块分散存放，其中任意 3 块可以同时消失};
  \foreach \x/\y in {0.45/2.21,1.45/1.48,2.45/0.99,3.45/0.74,4.45/0.72,5.45/0.95,6.45/1.42,7.45/2.13,8.45/3.08}{
    \draw[black!18] (\x,3.34) -- (\x,\y);
  }
  \draw[thick,blue!65] plot coordinates {(0.2,2.43) (0.7,2.01) (1.2,1.64) (1.7,1.33) (2.2,1.09) (2.7,0.90) (3.2,0.78) (3.7,0.71) (4.2,0.71) (4.7,0.76) (5.2,0.87) (5.7,1.05) (6.2,1.28) (6.7,1.57) (7.2,1.93) (7.7,2.34) (8.2,2.82) (8.7,3.35)};
  \foreach \x/\y in {0.45/2.21,1.45/1.48,3.45/0.74,4.45/0.72,6.45/1.42,7.45/2.13}{\fill[blue!75] (\x,\y) circle (2.2pt);}
  \foreach \x/\y in {2.45/0.99,5.45/0.95,8.45/3.08}{\draw[black!60,fill=white] (\x,\y) circle (2.4pt);}
\end{tikzpicture}

\emph{六个幸存的取值把那条次数小于 6 的曲线钉死，空心的三个是算回来的，不是存回来的。哪三块消失并不重要，重要的只是还剩几块——第一节那幅``任取 \(k\) 个点都行''的图，在这里就是运维承诺本身。}
\end{center}

\subsection{距离的天花板：Singleton 界}

先问一个与 Reed--Solomon 无关的问题：长度 \(n\)、消息长度 \(k\) 的码，最小距离最多能有多大？

\begin{theorem}[Singleton 界]
任何 \(q^k\) 个码字、长度为 \(n\) 的码，其最小距离满足 \(d\le n-k+1\)。
\end{theorem}

论证一行就够。把每个码字的最后 \(d-1\) 个坐标删掉，剩下的截断串必定两两不同——若有两个截断后相同，原来那两个码字至多在 \(d-1\) 个位置上不同，与最小距离是 \(d\) 冲突。于是 \(q^k\) 个互异的截断串挤在长度 \(n-d+1\) 的空间里，必有 \(q^k\le q^{\,n-d+1}\)，取对数即得。

这条界说的是一句朴素的话：每多买一分距离，就得多付一个冗余符号，且付出的绝不会比换来的少。达到等号的码称为 MDS 码（maximum distance separable），它把 \(n-k\) 个冗余符号一分不剩地换成了距离。

\subsection{多项式恰好顶到这个天花板}

Reed--Solomon 达到等号，理由上一节已经备好。取两个不同的消息多项式 \(f,g\)，它们对应的码字在哪些位置相同？正是差多项式 \(h=f-g\) 的根所在。\(h\) 非零且次数至多 \(k-1\)，因此至多有 \(k-1\) 个根，两个码字至少在 \(n-k+1\) 个位置上不同。与 Singleton 界一夹，
\[
d=n-k+1 .
\]

于是能力表可以直接照抄上一节的三条读法：检测至多 \(n-k\) 个错误，纠正至多 \(\lfloor(n-k)/2\rfloor\) 个错误，恢复至多 \(n-k\) 个擦除。若错误与擦除同时发生，判据合成一条预算规则
\[
2e+s\le n-k ,
\]
每个擦除花一个冗余符号，每个错误因为位置未知要花两个。开头那个 6 加 3 的条带里 \(n-k=3\)，三块盘掉线正好花光预算；但若某块盘悄悄返回了错数据而没有报错，一个这样的错误就要吃掉两个符号的余量。存储系统为此在块级另加校验和：与其让纠删码去猜哪块坏了，不如让它先自报家门，把一个错误降格成一个擦除，成本立刻减半。

\subsection{位置未知才是译码变难的地方}

只有擦除时，译码没有故事可讲——删掉缺失位置，对剩下的任意 \(k\) 个点做插值即可，这也是分布式存储里最常见的情形。

错误位置未知时，两个未知量搅在一起：不知道正确的多项式是什么，也不知道该信哪些位置。穷举所有 \(k\) 元子集去拟合当然可行，但组合数很快就不像话。代数译码的做法是把两个未知量绑成一组线性约束：设错误发生在集合 \(\mathcal E\) 对应的求值点上，构造错误定位多项式 \(E(x)=\prod_{i\in\mathcal E}(x-\alpha_i)\)，再令 \(Q=Ef\)。这样一来，无论某个位置正确与否，都有
\[
Q(\alpha_i)=E(\alpha_i)\,r_i ,
\]
正确的位置靠 \(r_i=f(\alpha_i)\) 成立，错误的位置靠 \(E(\alpha_i)=0\) 成立。未知的 \(E\) 与 \(Q\) 的系数满足一组线性方程，解出来再做一次除法 \(f=Q/E\)。这是 Berlekamp--Welch 的思路；工程上更快的伴随式译码、Euclid 算法、Berlekamp--Massey 各有各的省法，但都在利用同一个代数结构，没有谁在码字空间里盲目搜索。

顺带说一句边界：以上保证都是``错误数在预算内''时的确定性结论。超出预算，译码器可能给出一个错误的码字而毫无察觉，这是所有唯一译码的共同软肋，与算法好坏无关。

\subsection{它在哪些地方真的在跑}

Reed--Solomon 的符号通常取自 \(\mathbb F_{2^8}\)，一个符号一个字节。一串连续损坏的比特若集中在少数几个字节里，代价只按受损符号数计——这让它天然适合突发损坏，也就是二维码上的一块污渍、光盘上的一道划痕、网络里丢掉的整个分组。HDFS 与 Ceph 的纠删码、以及各种对象存储的跨机架冗余，用的都是它。

同一构造换个方向读，就是上一节末尾的 Shamir 秘密共享：那里希望不足 \(k\) 份时什么都得不到，这里希望凑够 \(k\) 份时什么都能还原，共用的是``\(k\) 个点定一条曲线''这一条。

它也有明确的用武之地边界。求值点必须互异，码长因此被 \(q\) 卡住，想要几万位的长码就得换更大的域，运算随之变贵；符号级的代数结构又使它不擅长利用软信息——信道给出的``这一位有七成把握是 0''这类置信度，在求根的框架里没有位置。长二进制信道上，用大量简单的局部约束换取廉价的迭代译码，往往比顶到 Singleton 界更划算。下一节就走这条路。
